%=======================================================================
\section{Discussion}
\label{sec:discussion}

\subsection{What Measuring the Paradigm Buys}

The trajectory paradigms of Section~\ref{sec:paradigms} are the same
vocabulary the literature already uses; the contribution is that they are
predicates over released per-frame quantities rather than adjectives attached
to a run. Three consequences follow, and all three appear in
Section~\ref{sec:results}.

First, intent becomes auditable. A run designed for co-observation measured
74\,\% and 77\,\% co-observation. That number cannot be produced by a
qualitative grading, and it is the number a user of the dataset needs in order
to know how much of a ``co-observation sequence'' is actually the condition
they wanted to train or test on.

Second, performance can be conditioned on the situation instead of averaged
over it. Because the strata are cut by $\omega$, $r$, $\beta$, $C$, and link
clearance rather than by the label, a method's accuracy under
$\beta > 90^\circ$ or beyond a node's evidence cap is directly reportable, and
does not inherit whatever mixture of conditions a run happened to contain. The
label is a description; the metric columns are the interface.

Third, negative and mixed results survive. Finding~1 says the geometry we
staged for is not the geometry that limited the team. A dataset that only
recorded its own design intent has no way to express that, and the effect would
have been reported as strong co-observation coverage.

\subsection{The Evidence-Definition Problem}

Finding~2 is, we think, the result with the widest reach beyond this dataset.
Collaboration rates---how often a partner supplies what the ego cannot---are
reported throughout the collaborative perception literature, and they are
usually stated as properties of a scene or a trajectory. They are not. A
$36\times$ swing between two modalities carried by the same robot on the same
run means such a rate is a joint property of the geometry \emph{and} the
sensing budget \emph{and} the evidence threshold, and that a rate quoted without
all three is not interpretable. Our response is to release the components
rather than only the rate: per (frame, node, object) evidence counts for every
modality, the empirically calibrated range cap for each, and the floor used.
Any rate in this paper can be recomputed under a different definition, and any
comparison against another dataset can be made at a matched definition.

The corresponding methodological point for range caps is that they must be
calibrated, not quoted from datasheets. Ours are set from the range at which
median in-box evidence falls to its floor and then frozen; the discrepancy
between a sensor's nominal range and the range at which it stops returning an
annotated chair is exactly the gap Finding~1 lives in.

\subsection{Threats to Validity}

We state the limits of the current release rather than defer them.

\paragraph{Sample size.} Two annotated static objects make the co-visibility
fraction $C$ a three-valued quantity ($0$, $0.5$, $1$), and every rate derived
from it is correspondingly coarse. The occlusion claims rest on 9 and 8
episodes, not on 733 frames; frames within an episode are not independent
samples, and we report episodes for that reason. One \SI{156}{\second} run at
one site is thin, and none of the numbers in Section~\ref{sec:results} should
be read as site- or platform-general.

\paragraph{Uncomputable conditions.} Doppler diversity is undefined here: the
scene contains no motion, so predicate~1 never fires, and \texttt{mobile\_2}
carries no radar and could not participate in the test even if it did.
Handover is under-sampled with a single infrastructure node, which also leaves
$\beta > 120^\circ$ thinly populated. Paradigms~2 and~5 are in a middle state:
the link clearance they depend on is computed, validated
(Table~\ref{tab:validation}), and released per frame, but it is not yet folded
into the per-frame classifier, so those labels do not appear in
Table~\ref{tab:seq} even though their ingredients do.

\paragraph{Map and calibration workarounds.} The occupancy grid is built from a
registered mobile-mapping cloud, and \SI{4.9}{\percent} of its points are
removed as statistical outliers with visible speckle remaining. The density
threshold, the consecutive-hit requirement, the \SI{0.6}{\metre} self-clearance
around each sensor origin, and the \SI{0.7}{\metre} clearance around both link
endpoints all exist because of map noise and because fixed nodes---the
infrastructure camera and the access point, the latter resting on a mapped
chair---are themselves mapped as structure. Each is a workaround with a proper
fix (a survey-grade scan; carving the map around known fixed nodes), and each
changes the reported numbers, which is why the parameters are frozen and
published with the data. Antenna extrinsics are currently approximated at the
camera origin, and per-frame throughput could not be measured: \texttt{iperf}
fired 10 times in \SI{156}{\second}, so goodput is not usable at frame rate and
we report RSSI and geometric clearance instead.

\paragraph{Ego symmetry.} Reporting both ego assignments of one bag controls
for platform but not for trajectory: the two nodes did not drive the same path.
The stronger control---paired repeats along one trajectory with a single factor
varied---is planned and not yet collected.

\subsection{What the Next Collection Must Add}

Four changes make every claim in the paradigm table computable, and they are
the priority ordering for the next campaign. \emph{Scripted moving
participants} (4--6 people, with one lane taped perpendicular to one
infrastructure boresight and along another) are required for Doppler diversity
to exist at all; it will not occur by accident, and it is the only condition in
Table~\ref{tab:paradigms} that no comparable dataset measures. \emph{Full
annotation of a declared bounded region} at 8--15 objects per frame turns $C$
from a coin flip into a distribution and makes every rate finer. A \emph{second
infrastructure node} makes handover testable, populates the wide-baseline tail
of $\beta$, and---by varying count at fixed elevation---separates the two
factors Finding~3 shows are currently confounded. \emph{Radar on
\texttt{mobile\_2}} lets both mobile nodes participate in any Doppler test.
Staging discipline improves the rest: marking blocking positions and logging
staged events on site should raise 6--8 deliberate occlusion events per
seven-minute run against this run's incidental 9, and paired repeats with one
varied factor (lighting, crowd density, node outage) turn the release from a
set of sequences into a set of controlled comparisons. Everything else on our
list---denser link sampling, logged antenna extrinsics, a survey-grade
map---improves the quality of claims we can already make rather than enabling
new ones.
